%% kapa: exam=1 round=29 year=2018 period=2 form=A subject=accounting subject_ko=회계학 seq=25 number=65
%% source: Q-Net 감정평가사 기출문제 (cst00309) artlSeq=5202808 file=extracted/1008/5202808/2018년 제29회 감정평가사 1차시험 기출문제/2018년 제29회 감정평가사 1차시험 기출문제/2교시 A형.pdf
%% conversion: kordoc 4.24.51 -> tools/gen1.py (auto); review=proofread-sonnet; flags=none
\begin{question}{65}
생물자산에 관한 설명으로 옳지 않은 것은?
\begin{choices}
  \item 생물자산의 순공정가치를 산정할 때에 추정 매각부대원가를 차감하기 때문에 생물자산의 최초 인식시점에 손실이 발생할 수 있다.
  \item 수확시점의 수확물은 어떠한 경우에도 순공정가치로 측정한다.
  \item 최초 인식후 생물자산의 순공정가치 변동으로 발생하는 평가손익은 발생한 기간의 당기손익에 반영한다.
  \item 순공정가치로 측정하는 생물자산과 관련된 정부보조금에 다른 조건이 없는 경우에는 이를 수취할 수 있게 되는 시점에 기타포괄손익으로 인식한다.
  \item 생물자산을 최초로 인식하는 시점에 시장 공시가격을 구할 수 없고, 대체적인 공정가치측정치가 명백히 신뢰성 없게 결정되는 경우에는 원가에서 감가상각누계액과 손상차손누계액을 차감한 금액으로 측정한다.
\end{choices}
\end{question}
\answer{4}
